\documentclass[12pt]{article}
\usepackage{geometry}
\geometry{margin=1.4in}
\usepackage{setspace}
\setstretch{1.4}

\title{Energy Systems as Displacement Traps:\\
Why Renewables Cannot Replace Fossil Fuels Without Structural Change}
\author{Diego Rincón\\
\small Independent}
\date{}

\begin{document}

\maketitle

\begin{abstract}
Energy systems displace toward fossil fuels not because they are better but because they are cheaper to maintain. Coal, oil, and gas are locked-in ground states: the infrastructure is built, the supply chains are established, the cost of operations is low. Renewable energy is the intended ground state $S_0$, but the return cost from the fossil ground state $S^*$ is prohibitive. The transition is not a technical problem. It is a displacement problem.
\end{abstract}

\section{The Cheap Ground State}

Fossil fuels created a massive stable state:
- Extraction infrastructure (mines, wells, refineries)
- Transportation networks (pipelines, tankers, rail)
- Power plants (coal, gas)
- Consumer expectations (cheap electricity, cheap gas)
- Political alignment (subsidies, regulatory capture)

This ground state is cheap to maintain. Each day, staying in it costs less than leaving.

\section{The Expensive Return}

Renewable energy requires:
- New infrastructure (solar, wind, storage)
- New grid systems (distributed, smart)
- New extraction sites (lithium, rare earths)
- New consumer expectations (variable supply, higher cost)
- Political realignment (remove fossil subsidies)

The cost to return to $S_0$ is enormous. It exceeds the annual climate budget we're willing to spend.

\section{Why Net-Zero Fails}

Pledges to reach net-zero by 2050 assume a linear path. But systems don't move linearly. They displace in jumps. The system will reach a new stable state somewhere between fossil fuels and renewables: hybrid, mixed, contradictory.

This is not failure. This is the system finding an attractor with lower return cost than pure renewables.

\section{The Only Path}

Returning to $S_0$ requires:
\begin{enumerate}
\item Removing all subsidies from fossil fuels
\item Pricing carbon at true cost (not politically acceptable cost)
\item Building new infrastructure faster than the old decays
\item Accepting temporary energy poverty in transition regions
\item International coordination (no cheap coal from elsewhere)
\end{enumerate}

Each step raises the cost of return. The system will stay displaced unless forced by crisis or cost inversion (renewables cheaper than fossils for enough applications that the entire infrastructure tips).

\end{document}
